
Uncertainty-based hallucination detectors achieve strong in-distribution performance but lack principled criteria for predicting cross-benchmark transfer success. This paper presents a systematic investigation of this problem, discovering that benchmarks cluster into families based on uncertainty distribution similarity and that transfer success depends on cluster membership. Using semantic entropy and Jensen-Shannon divergence, six QA benchmarks were clustered into two families---Factual Recall (TriviaQA, Natural Questions, SQuAD) and Entity/Claim Verification (PopQA, HaluEval-QA, FEVER)---with silhouette score 0.8245. Within-cluster threshold transfer showed mean AUROC degradation of 0.032, while cross-cluster transfer exhibited degradation of 0.223---a ratio of approximately $7\times$. These findings suggest that cluster membership, determined via JS-divergence, may provide a practical criterion for deployment decisions: thresholds can be reused within clusters but require recalibration when crossing cluster boundaries.

